\uuid{Qs4I}
\exo7id{3527}
\titre{exo7 3527}
\auteur{quercia}
\organisation{exo7}
\datecreate{2010-03-10}
\isIndication{false}
\isCorrection{true}
\chapitre{Réduction d'endomorphisme, polynôme annulateur}
\sousChapitre{Polynôme caractéristique, théorème de Cayley-Hamilton}
\cpge{mpsi-mp2i.matrices-applications-lineaires.b, pcsi.matrices-applications-lineaires.b, ptsi.matrices-applications-lineaires.b, mpsi-mp2i.changements-bases.c, bcpst1.matrices, mp-mpi.reduction-endomorphismes.b, pc.reduction.a, psi.reduction-endomorphismes.a, pt.reduction-endomorphismes.a, mp-mpi.reduction-endomorphismes.d, pc.reduction.c, psi.reduction-endomorphismes.c, pt.reduction-endomorphismes.c, bcpst2.valeurs-propres-vecteurs-propres.a, bcpst2.valeurs-propres-vecteurs-propres.b}
\module{Algèbre}
\niveau{L2}
\difficulte{}

\contenu{
\question{
Soit $A=(a_{ij}) \in \mathcal{M}_n(\R)$ telle que $a_{ij} = \frac ij$.
    $A$ est-elle diagonalisable~?
}
\reponse{
$\mathrm{rg} A = 1$ donc $\dim\mathrm{Ker} A = n-1$ et $0$ est valeur propre
             d'ordre au moins $n-1$. La somme des valeurs propres est
             $\mathrm{tr} A = n$ donc la drenière valeur propre est $n$ et le
             sous-espace propre associé est de dimension~$1$. Donc $A$
             est diagonalisable.
}
}
